\subsection{Sockets TCP}

Para implementar la comunicación mediante sockets TCP se
desarrolló un servidor y un cliente en Python, utilizando
el módulo \texttt{socket}, aqui cada nodo se ejecuta como un proceso
independiente y dispone de ambos componentes lo que permite
enviar y recibir mensajes durante las pruebas.

\subsubsection{Servidor TCP}

El servidor se encarga de escuchar las conexiones entrantes
en la dirección \texttt{127.0.0.1} y en el puerto asignado
a cada nodo. Cuando recibe una conexión crea un hilo para
atenderla de modo que puede gestionar varias conexiones
sin detener la aceptación de nuevos clientes.

Los mensajes se reciben en formato JSON y separados por
saltos de línea. Antes de procesarlos el servidor comprueba
que contengan los campos \texttt{sender}, \texttt{timestamp}
y \texttt{message}, con los tipos de datos esperados.

Si el mensaje es válido, se entrega a la lógica común
del nodo y se devuelve una confirmación con el estado
\texttt{ok}. Cuando se detecta un formato incorrecto
el servidor responde con un estado de error.

\subsubsection{Cliente TCP}

El cliente establece una conexión TCP con el nodo de destino
y la conserva para reutilizarla en los siguientes envíos.
Esta decisión evita crear una nueva conexión por cada mensaje
y permite mantener el mismo criterio durante las pruebas
de comunicación.

Cada mensaje contiene el identificador del emisor, la marca
de tiempo lógica de Lamport y el contenido que se desea
transmitir. Después del envío el cliente espera la
confirmación del receptor y verifica que corresponda
al destino esperado.

Si la conexión inicial falla, el programa utiliza el
mecanismo de reintentos definido en la lógica común.
Al finalizar la ejecución, se cierran las conexiones
que permanecen abiertas.

\subsubsection{Integración con el reloj de Lamport}

La implementación TCP utiliza la clase \texttt{NodoComun},
la cual centraliza las operaciones compartidas con gRPC.
Antes de enviar un mensaje, esta clase actualiza el reloj
lógico y proporciona la marca de tiempo que debe
transmitirse.

Cuando se recibe un mensaje su marca de tiempo se utiliza
para actualizar el reloj local. En el modo de intercambio
Lamport, los mensajes entrantes se almacenan primero en
una cola y posteriormente son procesados por el hilo
principal del nodo, respetando el orden de llegada.

Las confirmaciones TCP se utilizan para comprobar la
recepción de los mensajes pero no se registran como
eventos adicionales del reloj de Lamport.

\subsubsection{Registro de latencia}

Para las pruebas de rendimiento se mide el tiempo de ida
y vuelta (\textit{Round-Trip Time}, RTT) desde el envío
del mensaje hasta la recepción de su confirmación.

La medición se realiza mediante
\texttt{time.perf\_counter()}, y los resultados se
almacenan en archivos CSV. Cada registro incluye
información del mecanismo, escenario, tamaño del mensaje,
corrida, nodo, número de envío y latencia en milisegundos.

Este procedimiento permite comparar los resultados
obtenidos mediante sockets TCP con los de gRPC
utilizando la misma lógica experimental.